\scanpage{76}

% L01 -- top, parenthesis continued
(L'indétermination est donc section d'un rev. principal de $S$ de \unsure{gpe|groupe} $\mu_2$, défini mod él. de $\mu_2(S)$ \ldots). Donc $s_0$ est alors donné par la formule $(*)$.

% L02 -- upper third
(Rien n'empêche d'ailleurs que $f_V(s_0)=0$)

% L03 -- upper middle, struck heading
\deleted{(E) Supposons $n$ impair}

% L04 -- upper middle
Si on se donne un hom

% L05 -- upper middle, homomorphism
\[
\varphi_G^C:G\unsure{-|\longrightarrow}\mathrm{Aut}(C,Q_C)
\]

% L06 -- middle, square brackets
[$C$ fibré proj. de rg $n$, $Q_C$ \inserted{hyper}quadrique,]

% L07 -- above list
satisfaisant les conditions

% L08 -- first condition
1) Les $\tau_i^X$ réflexions \deleted{\unsure{aut|orth.}} \qquad ($0\le\unsure{j|i}\le n$) -- d'où $\ell_i\in\Gamma C,\ H_i$ hyperplans)

% L09 -- second condition
2) les $\ell_i$ engendrent $C$, \deleted{\unsure{les|i.e.}} $\bigcap_{0\le i\le n}H_i^{\unsure{C}}=\varnothing$

% L10 -- third condition
3) \deleted{\unsure{$\ell_i$}} $\ell_i\notin H_{i+1}$ \qquad $0\le i\le n-1$

% L11 -- fourth condition, crossed out
\deleted{\unsure{et} 4) \deleted{\unsure{un}} $\ell_0\notin Q_C$ en chaque fibre}.

% L12 -- lower left, diagonal qualification
% slanted
\margin{Pour $n\ge2$ (si $n=1$, une réflexion de la droite \unsure{pr.} $C=\mathrm P(V)$ ne se remonte pas de \uline{façon univoque} en une réfl. de $V$ \ldots]}

% L13 -- lower part
alors loc. (\unsure{ét|et}) ça provient de $(V,f_V,(\tau_i^V))$ satisfaisant les conditions de (D), $V$ déterm. \unsure{près de|mod} $\otimes L$ ($L$ inv), $f_V$ défini mod unité. On \unsure{Utilisant|utilise} \unsure{la remarque|la remarque de|alors} ; $L_0\subset V$, on \unsure{les|l'} ``\unsure{épingler|épingle|épingles}'' i.e. éliminer l'indét. $\otimes L$, en prenant $L\simeq\mathcal O_S$ i.e. section $a_0$ de $L_0$.
% archive sheet number: 75
